\section{The same lower threshold for canonical finite priors}
\label{sec:canonical-negative}
The negative theorem also applies at the actual canonical scales used by the finite positive theorem. This removes the unspecified small-scale threshold in those regimes.

\begin{theorem}[Finite-scale failure at every temperature]
\label{thm:canonical-negative}
Use the law and the sufficiently large sample counts of Theorem~\ref{thm:lower}. With the same sample probability, simultaneously for every $N\ge N_*(n)$, every $\gamma\ge2$, and every $\lambda>0$, the posterior at $\sigma^2=N^{-\gamma}$ has all the stated raw and clipped risks at least $p+1/100$. When $H=1$, the assertion holds for every $\gamma\ge1$.

No prior-strength or width limit is taken. The temperature, width, and scaling exponent may be selected from the sample, provided they lie in these ranges.
\end{theorem}
\begin{proof}
Put $D_0=d+H+1$. The regimes in the theorem satisfy $P_N\sigma^2\le D_0$ and $N\sigma^2\le1$.

For the low-temperature branch, use the exact balanced memorizer of energy at most $K_{\mathrm{mem}}$. Its empirical loss is at most one: fitting the noisy atom means reduces the zero predictor's noisy loss, and it vanishes on the steep component. Lemma~\ref{lem:finite-perturbation} gives the output root mean square perturbation bound $rK_m$, where
\[
 K_m=\sqrt2 L(\sqrt{K_{\mathrm{mem}}/L}+1)^H.
\]
Take the full-dimensional Gaussian comparison measure centered at this memorizer, with covariance $r_n^2$ times the prior covariance, where
\[
 r_n=\frac1{8n(K_m+1)}.
\]
Minkowski bounds its expected empirical loss by the memorizer's loss plus $1/n$. The exact Gaussian relative entropy and Gibbs variational identity therefore give
\[
 \int\empRisk_S\,d\Pi\le
 \empRisk_S(z_{w_{\mathrm{mem}}})+\frac1n
   +\lambda\left[K_{\mathrm{mem}}+2D_0\log(1/r_n)\right].
\]
Dropping the within-atom residual $W/n$, this bounds the posterior expected noisy residual by
\[
 \frac{n_s}n+\frac1n+
       \tau_n\left[K_{\mathrm{mem}}+2D_0\log(1/r_n)\right]
       =q+o(1)
 \quad(\lambda\le\tau_n).
\]
Here $K_m$ grows polynomially with $n$ at fixed depth, so $\log(1/r_n)=O_H(\log n)$; the estimates are uniform over all widths and exponents in the theorem. For large $n$ the residual is at most $1/5$ and $S_1/n\ge1/3$. Jensen and the reverse-triangle proof above give the same excess lower bound $p(\sqrt{1/3}-\sqrt{1/5})^2>1/100$ for all stated prediction conventions.

For large temperatures, prior expected empirical loss is at most three, conservatively, by the second-moment calculation in Section~\ref{sec:finite-positive}. Jensen gives $Z\ge e^{-3/(2\lambda\sigma^2)}$. Combining this with the Gaussian energy moment yields
\[
 \int e^{Q_N/(4\sigma^2)}\,d\Pi
      \le2^{P_N/2}e^{3/(2\lambda\sigma^2)}.
\]
Thus the posterior law of $Q_N$ is stochastically dominated by
\[
 2P_N\sigma^2\log2+\frac6\lambda+V_\sigma,
 \qquad V_\sigma\sim\operatorname{Exp}(\text{mean }4\sigma^2).
\]
For $j=\lceil L/2\rceil$, Minkowski and
$\|\operatorname{Exp}(\text{mean }1)\|_{L/2}\le(j!)^{1/j}\le j$
give
\[
 \left(\int Q_N^{L/2}\,d\Pi\right)^{2/L}
    \le2D_0\log2+\frac6\lambda+4j.
\]
For $\lambda\ge\tau_n$, define $B_n=2D_0\log2+6/\tau_n+4j$.
The input-coordinate sensitivity bound shows that the half-difference of the raw posterior mean at each inner pair is at most
\[
 \eta_{\mathrm{fin}}=\delta d^{-1/2}(B_n/L)^{L/2}\longrightarrow0.
\]
This follows from the same strict exponent condition as before, since $1/\tau_n$ dominates $D_0$ for large $n$. The half-difference bound also applies to the mean of clipped outputs and to clipping the raw mean. For large $n$ it is at most $1/2$, giving excess at least $q/8>1/100$. Jensen then gives the corresponding Gibbs-risk bounds. The argument is uniform over every parameter in the theorem, and proves the claim.
\end{proof}
